\bookchapter{第六十五章：專案一——我的記帳本（本地版）}

\chapterintro{歡迎來到你的第一個“正經專案”。這個記帳本雖然不聯網，但五臟俱全：能記一筆賬、能看流水、能按月統計、能刪改。做完它，你就完整走了一遍“資料 邏輯 介面”的全流程。別急著抄程式碼——先想清楚再動手，這是咱們一貫的原則。}

\booksection{一、需求思考：先問五個問題}

\begin{center}
\begin{tabularx}{\linewidth}{|X|X|}
\hline
\textbf{要素} & \textbf{答案} \tabularnewline
\hline
Who（給誰用） & 自己用，家人朋友也行 \tabularnewline
What（做什麼） & 記收支、看流水、按月彙總、按分類統計 \tabularnewline
Why（為什麼） & 練手 + 搞清楚錢花哪兒了 \tabularnewline
How（怎麼做） & Swift + SwiftUI + SwiftData + MVVM，純本地無後端 \tabularnewline
Constraint（約束） & 一週內完成；最低支援 iOS 17（SwiftData 門檻） \tabularnewline
\hline
\end{tabularx}
\end{center}

\begin{center}\small\color{gray}記帳本需求五要素\end{center}

然後拆解成“每一步都能用”的小步驟：(1) 建 SwiftData 資料層（能存一筆賬）；(2) 列表頁（能看流水）；(3) 錄入跳出視窗（能加賬）；(4) 刪除與編輯；(5) 月度彙總卡片。每完成一步，App 都是一個“能用的東西”。

\booksection{二、資料層：讓“賬”有個家}

先定義一筆賬長什麼樣——欄位要夠用但別貪多：

\begin{lstlisting}[language=Swift]
import SwiftData

@Model
class Record {
    var type: String      // "income" 收入 | "expense" 支出
    var amount: Double    // 金額（元）
    var category: String  // 分類：餐飲、交通、工資……
    var note: String      // 備註（可選）
    var date: Date        // 記帳日期

    init(type: String, amount: Double, category: String, note: String = "", date: Date = Date()) {
        self.type = type
        self.amount = amount
        self.category = category
        self.note = note
        self.date = date
    }
}
\end{lstlisting}

\begin{tcbinfo}{小貼士：金額為什麼用 Double？}
嚴謹的金融 App 會用 Decimal（精確十進位制）來避免浮點誤差（0.1+0.2$\neq$0.3 的老問題）。練手專案用 Double 足夠；如果哪天做真金融產品，記得換成 Decimal 或“以分為單位的 Int”——這是無數前輩踩出來的坑。

\end{tcbinfo}

\booksection{三、MVVM：介面與邏輯“分家”}

直接讓介面運算元據庫？能跑，但介面程式碼會又臭又長。MVVM 架構讓各司其職：View（介面）只管畫，ViewModel（邏輯中樞）管資料，Model（SwiftData）管儲存。

View（SwiftUI 介面）

ViewModel（@Observable 邏輯中樞）

Model（SwiftData 資料層）

\begin{center}\small\color{gray}圖 4：MVVM 三層結構——資料單向流動，各層互不越界\end{center}

\begin{lstlisting}[language=Swift]
import Observation

@Observable
class RecordViewModel {
    var records: [Record] = []
    private let context: ModelContext

    init(context: ModelContext) { self.context = context }

    func load() {
        // SwiftData 的 #Predicate 篩選 + 排序
        let descriptor = FetchDescriptor<Record>(
            sortBy: [SortDescriptor(\.date, order: .reverse)]
        )
        records = (try? context.fetch(descriptor)) ?? []
    }

    func add(type: String, amount: Double, category: String, note: String) {
        let r = Record(type: type, amount: amount, category: category, note: note)
        context.insert(r)
        try? context.save()
        load()                       // 重新載入，介面重新整理
    }

    func delete(_ record: Record) {
        context.delete(record)
        try? context.save()
        load()
    }
}
\end{lstlisting}

介面裡 @State private var vm = RecordViewModel(...)，然後直接 vm.records 渲染列表——資料一變，介面自動重新整理。

\booksection{四、介面：列表頁 + 錄入跳出視窗}

\begin{lstlisting}[language=Swift]
struct RecordListView: View {
    @State private var vm: RecordViewModel
    @State private var showAdd = false

    var body: some View {
        NavigationStack {
            List {
                Section("本月結餘") {          // 彙總卡片
                    SummaryRow(records: vm.records)
                }
                Section("流水") {
                    ForEach(vm.records) { r in
                        RecordRow(record: r)
                    }
                    .onDelete { idx in
                        idx.forEach { vm.delete(vm.records[$0]) }
                    }
                }
            }
            .navigationTitle("我的記帳本")
            .toolbar {
                Button { showAdd = true } label: {
                    Image(systemName: "plus")
                }
            }
            .sheet(isPresented: $showAdd) {
                AddRecordView(vm: vm)          // 錄入彈窗
            }
        }
    }
}
\end{lstlisting}

錄入跳出視窗要點：金額輸入框鍵盤用 .keyboardType(.decimalPad)（數字鍵盤），收入/支出用分段控制元件（Picker），分類用選單選擇，點“儲存”調 vm.add(...) 並 dismiss() 關窗。

\begin{tcbwarn}{注意：輸入校驗別偷懶}
金額輸入框一定要校驗：空值、負數、“abc”這類非法輸入都要攔住（用 Double(text) 判斷，失敗返回 nil）。App 崩掉最常見的原因之一就是“把使用者的輸入當成了乖孩子”。記住：使用者輸入的一切，都要當敵人審查。

\end{tcbwarn}

\booksection{五、編輯、刪除與彙總}

\begin{itemize}
  \item 刪除：List 的 .onDelete 左滑刪除 + 確認跳出視窗（誤刪是記帳類 App 使用者最痛恨的事）；
  \item 編輯：點行進編輯跳出視窗，預填原值（改一個字比重新輸入省事）；
  \item 彙總：用 filter + reduce 算"本月收入 / 支出 / 結餘"，頂部卡片展示：
\end{itemize}

\begin{lstlisting}[language=Swift]
func monthSummary(_ records: [Record]) -> (income: Double, expense: Double) {
    let cal = Calendar.current
    let monthRecords = records.filter {
        cal.isDate($0.date, equalTo: Date(), toGranularity: .month)
    }
    let income = monthRecords
        .filter { $0.type == "income" }
        .reduce(0) { $0 + $1.amount }
    let expense = monthRecords
        .filter { $0.type == "expense" }
        .reduce(0) { $0 + $1.amount }
    return (income, expense)
}
\end{lstlisting}

\booksection{六、驗收：三層檢查 + 兩個進階挑戰}

專案做完，按“功能 邊界 架構”三層驗收：能不能記、能不能刪、重啟 App 資料還在不在（功能）；金額輸 0 或負數會不會崩（邊界）；介面和資料有沒有混在一起、欄位有沒有硬編碼（架構）。

基礎版驗收透過後，給自己加兩道“進階菜”：

\begin{enumerate}
  \item 分類統計：加一個"統計"標籤頁，用 Swift Charts 把分類支出畫成餅圖/柱狀圖——蘋果官方圖表庫，幾行程式碼出圖；
  \item iCloud 同步：SwiftData 開 iCloud 同步只需配幾項設定——帳本在 iPhone/iPad 之間自動同步，瞬間"進階"起來。
\end{enumerate}

\begin{tcbpractice}{來動手：本週完成它}
按“資料層 ViewModel 列表頁 錄入跳出視窗 彙總卡片”的順序動手，每天一步。卡住了就把報錯原樣丟給搜尋引擎或 AI（還記得第七篇《AI程式設計實戰篇》嗎？描述問題時說清“期望行為、實際行為、復現步驟”）。完成後用它記一週賬——自己做的 App 用起來，成就感會推著你繼續往前走。

\end{tcbpractice}

\begin{tcbsummary}{本章小結}
\begin{itemize}
  \item 需求五要素 拆成"每步能用"的小步驟，再動手
  \item 資料層：@Model 模型 + modelContainer 容器 + FetchDescriptor 查詢
  \item MVVM 分工：View 只管畫，@Observable 的 ViewModel 管邏輯，SwiftData 管儲存
  \item 介面：List + Section + toolbar + sheet 錄入跳出視窗
  \item 三層驗收（功能/邊界/架構）+ 進階挑戰（Swift Charts、iCloud 同步）
\end{itemize}

\end{tcbsummary}
